% generated by roam.py handbook export -- не редактировать, править заметки
% 20260922143839-дифференциальныи_критерии_сильнои_выпуклости.org (body 8cb481c501c7)

\begin{theorem}[{Дифференциальный критерий сильной выпуклости}]\label[theorem]{note:1B5A3E25-56E8-4D5E-9CF6-E2731DA77604}
  Пусть $E \subset \R^n$~--- \hyperref[note:316B5BFA-8B9F-4267-9086-D6593AAB5F9B]{выпуклое множество}. Если $f: E \rightarrow \R$~--- \hyperref[note:957FA09D-8883-40A5-B4C7-237D9F815A42]{дифференцируемая}, то
  $f$~--- $\mu$-\hyperref[note:22E546FC-186F-4A67-85AA-2C85DCEF9021]{сильно выпуклая} $\Leftrightarrow \forall x, y \in E\ f(y) \ge f(x) + \left\langle \nabla f(x); y-x \right\rangle + \frac{\mu}{2} \|y-x\|^2$
\end{theorem}

\begin{remark}
  Если для \hyperref[note:1CA14F2C-637F-41DE-9E97-72075D19E97F]{обычной выпуклости} мы подпирали функцию прямой (касательной), то в этом случае мы подпираем уже параболой
\end{remark}

\begin{remark}
  Если $\|y\| \rightarrow +\infty$, то правая часть тоже $\rightarrow +\infty$
\end{remark}

\begin{proof}
  Возьмем функцию $\phi(x) = f(x) - \frac{\mu}{2} \|x\|^2$. Она  дифференцируема.
  Докажем, что она выпуклая, используем \hyperref[note:18BD8497-CEBD-4962-8144-474B489606C3]{дифференциальный критерий выпуклости}.

  $\phi$~--- выпуклая $\Leftrightarrow \forall x,y \in E\ \phi(y) \ge \phi(x) + \left\langle \nabla \phi(x); y-x \right\rangle$

  Заметим $\nabla \phi(x) = \nabla f(x) - \mu x$, так как \hyperref[note:1756907B-1FF2-4F62-8553-759B8509281C]{градиент}~--- линеен, а градиент $\|x\|^2$~--- это $2x$.
  Давайте теперь раскроем $\phi$ в неравенстве:
  \[
  f(y) - \frac{\mu}{2} \| y \|^2 \ge f(x) - \frac{\mu}{2} \| x \|^2 + \left\langle \nabla f(x) - \mu x; y-x \right\rangle
  \]

  Давайте раскроем \hyperref[note:F2942E98-6115-4600-AA3F-E0F687D9C732]{скалярное произведение} в правой части:
  \[
  \left\langle \nabla f(x) - \mu x; y-x \right\rangle =  \left\langle \nabla f(x); y-x \right\rangle - \mu \left\langle x; y - x \right\rangle =  \left\langle \nabla f(x); y-x \right\rangle - \mu \left\langle x; y \right\rangle  +  \mu \left\langle x; x \right\rangle
  \]

  Тогда
  \[
  f(y) - \frac{\mu}{2} \| y \|^2  \ge f(x) - \frac{\mu}{2} \|x\|^2 + \left\langle \nabla f(x); y-x \right\rangle - \mu \left\langle x; y \right\rangle  +  \mu \left\langle x; x \right\rangle
  \]
  Перенесем норму $y$ на право и заменим $\mu \left\langle x;x \right\rangle = \mu \|x \|^{2}$,
  \[
  f(y) \ge f(x) - \frac{\mu}{2} \| x \|^2 + \mu \| x \|^2 + \frac{\mu}{2} \| y \|^2 + \left\langle \nabla f(x); y-x \right\rangle - \mu \left\langle x;y \right\rangle
  \]
  Итого получим
  \[
  f (y) \ge f(x) + \frac{\mu}{2} \| x \|^2 + \frac{\mu}{2} \| y \|^2 - \mu \left\langle x;y \right\rangle + \left\langle \nabla f(x); y-x \right\rangle
  \]

  Еще заметим, что  $\|y-x\|^2 = \|x\|^2 +\|y\|^2 - 2 \left\langle x;y \right\rangle$, давайте используем эту формулу \\
  Итого
  \[
  f(y) \ge f(x) + \frac{\mu}{2} \|y-x\|^2 + \left\langle \nabla f(x); y-x \right\rangle
  \]
  Все переходы~--- равносильности, поэтому получили в точности неравенство из условия: оба направления доказаны
\end{proof}
